\nocite{8b5861fc}
\nocite{c7f15065}
In this section we recount the mathematics needed to obtain an impression of topological insulators.
The mathematical branches relevant for the theory of topological insulators are
\begin{enumerate}
\item[$\bullet$]
  (functional) analysis
\item[$\bullet$]
  homotopy theory
\end{enumerate}
The former is relevant for describing the quantum-mechanical aspects and the latter for the topological ones.
And while we assume the reader to be familiar with functional analysis (in particular spectral theory and the mathematics of quantum mechanics) and basic homotopy theory (see e.g. \cite{catctxphy} or the standard reference \cite{8b5861fc}), we do not require familiarity with more advanced homotopy theory.
Most notably, we do not assume any knowledge of topological K-theory.
Topological K-theory is the main tool for a rough understanding of topological phases, and it is the subject of this section.
Our presentation is essentially an abstraction of \cite{c7f15065}.
Accordingly, we include subsections on categories, the Grothendieck group, and vector bundles, leading up to topological K-theory in the final subsection.
Note that our treatment of these topics is very high-level and mostly algebraic, largely glossing over geometric aspects.
This suffices for our purposes, but the reader may eventually wish to consult a source such as \cite{c7f15065} for a deeper understanding.
In fact, characteristic classes - in particular Chern classes - as presented there are important for a thorough understanding of our account of topological insulators in section \ref{sec:topins}.
This also involves some differential geometry, but for a high-level picture both topics can be neglected.
\\
\begin{rem}
\label{rem:style}
Before getting started, a few words about the mathematical style: 
\begin{enumerate}
\item[$\bullet$]
  Mathematical definitions use \textbf{boldface type} to indicate the term being defined.
In later sections we will sometimes encounter new terms that we cannot define here for various reasons.
In this case we use \textit{italic type} instead of boldface type and refer the reader to other sources.
In these preliminaries we usually provide an example for each definition for the sake of comprehension.
However, in most cases we do not give proofs of the stated propositions.
If anything, we only sketch them. 
\item[$\bullet$]
  $\mathsf{A}$ denotes both finite and infinite index sets, that is, arbitrary sets with specified cardinality.
$\mathbb{K}$ denotes either the real numbers $\mathbb{R}$ or the complex numbers $\mathbb{C}$, and $(\cdot \vert \cdot)$ denotes the scalar product.
Following common convention, we use the term \textit{space} for a set together with some \textit{structure} (e.g. a topology, vector space structure, \ldots), and \textit{subspace} means a subset equipped with the induced structure.
In this context, we use the term \textbf{map} to mean a structure-preserving function for the sake of brevity.
That is, in the context of topological spaces, a map means a continuous function, for example.
\end{enumerate}
\phantom{proven}
\hfill
$\square$
\end{rem}
